\documentclass{amsart}
% For dealing with references we use the comment environment
\usepackage{verbatim}
\newenvironment{reference}{\comment}{\endcomment}
%\newenvironment{reference}{}{}
\newenvironment{slogan}{\comment}{\endcomment}
\newenvironment{history}{\comment}{\endcomment}

% For commutative diagrams we use Xy-pic
\usepackage[all]{xy}

% We use 2cell for 2-commutative diagrams.
\xyoption{2cell}
\UseAllTwocells

% We use multicol for the list of chapters between chapters
\usepackage{multicol}

% This is generally recommended for better output
\usepackage{lmodern}
\usepackage[T1]{fontenc}

% For cross-file-references

% Package for hypertext links:
\usepackage{hyperref}

% For any local file, say "hello.tex" you want to link to please

% Theorem environments.
%
\theoremstyle{plain}
\newtheorem{theorem}[subsection]{Theorem}
\newtheorem{proposition}[subsection]{Proposition}
\newtheorem{lemma}[subsection]{Lemma}

\theoremstyle{definition}
\newtheorem{definition}[subsection]{Definition}
\newtheorem{example}[subsection]{Example}
\newtheorem{exercise}[subsection]{Exercise}
\newtheorem{situation}[subsection]{Situation}

\theoremstyle{remark}
\newtheorem{remark}[subsection]{Remark}
\newtheorem{remarks}[subsection]{Remarks}

\numberwithin{equation}{subsection}

% Macros
%
\def\lim{\mathop{\mathrm{lim}}\nolimits}
\def\colim{\mathop{\mathrm{colim}}\nolimits}
\def\Spec{\mathop{\mathrm{Spec}}}
\def\Hom{\mathop{\mathrm{Hom}}\nolimits}
\def\Ext{\mathop{\mathrm{Ext}}\nolimits}
\def\SheafHom{\mathop{\mathcal{H}\!\mathit{om}}\nolimits}
\def\SheafExt{\mathop{\mathcal{E}\!\mathit{xt}}\nolimits}
\def\Sch{\mathit{Sch}}
\def\Mor{\mathop{\mathrm{Mor}}\nolimits}
\def\Ob{\mathop{\mathrm{Ob}}\nolimits}
\def\Sh{\mathop{\mathit{Sh}}\nolimits}
\def\NL{\mathop{N\!L}\nolimits}
\def\CH{\mathop{\mathrm{CH}}\nolimits}
\def\proetale{{pro\text{-}\acute{e}tale}}
\def\etale{{\acute{e}tale}}
\def\QCoh{\mathit{QCoh}}
\def\Ker{\mathop{\mathrm{Ker}}}
\def\Im{\mathop{\mathrm{Im}}}
\def\Coker{\mathop{\mathrm{Coker}}}
\def\Coim{\mathop{\mathrm{Coim}}}

% Boxtimes
%
\DeclareMathSymbol{\boxtimes}{\mathbin}{AMSa}{"02}

%
% Macros for moduli stacks/spaces
%
\def\QCohstack{\mathcal{QC}\!\mathit{oh}}
\def\Cohstack{\mathcal{C}\!\mathit{oh}}
\def\Spacesstack{\mathcal{S}\!\mathit{paces}}
\def\Quotfunctor{\mathrm{Quot}}
\def\Hilbfunctor{\mathrm{Hilb}}
\def\Curvesstack{\mathcal{C}\!\mathit{urves}}
\def\Polarizedstack{\mathcal{P}\!\mathit{olarized}}
\def\Complexesstack{\mathcal{C}\!\mathit{omplexes}}
% \Pic is the operator that assigns to X its picard group, usage \Pic(X)
% \Picardstack_{X/B} denotes the Picard stack of X over B
% \Picardfunctor_{X/B} denotes the Picard functor of X over B
\def\Pic{\mathop{\mathrm{Pic}}\nolimits}
\def\Picardstack{\mathcal{P}\!\mathit{ic}}
\def\Picardfunctor{\mathrm{Pic}}
\def\Deformationcategory{\mathcal{D}\!\mathit{ef}}

\setlength{\emergencystretch}{3em}
\begin{document}
\title[Stacks Project, Tag 0E9L]{250 --- Expression of algebraic upper shriek by dualizing biduality}
\maketitle
\noindent Copyright (C) 2005--2025 Johan de Jong. Stacks Project authors. Source: \href{https://stacks.math.columbia.edu/tag/0E9L}{Tag 0E9L}.

\noindent Editorial difficulty note (not author text). Difficulty H2: The polynomial-algebra case is proved through bounded injective models and a truncation argument valid for potentially unbounded coherent complexes. The surjection case uses a derived tensor-Hom adjunction chain and Yoneda; composition then assembles arbitrary finite-type maps. Establishing this functorial identification demands more than the existence of dualizing complexes..

\noindent Editorial provenance note (not author text). History: excerpt from revision \texttt{a0444\allowbreak 6e57e\allowbreak c1fbc\allowbreak 252a8\allowbreak 71afc\allowbreak ec775\allowbreak 2fb28\allowbreak 07b14}, dualizing.tex, lines 4693--4783. Reference presentation was adapted; original-author context and core support blocks were retained or added. The mathematical wording is unchanged. Licensed under GNU FDL 1.2 or later; no invariant sections or cover texts. See ../licenses/COPYING. Context blocks reproduce author definitions and supporting results; they are not additional counted problems.

\section{Upper shriek algebraically}
\label{section-relative-dualizing-complex-algebraic}

\noindent
For a finite type homomorphism $R \to A$ of Noetherian rings
we will construct a functor $\varphi^! : D(R) \to D(A)$
well defined up to nonunique isomorphism which
as we will see in Duality for Schemes, Remark
\href{https://stacks.math.columbia.edu/tag/0BV2}{remark local calculation shriek (Tag 0BV2)}
agrees up to isomorphism with the upper shriek functors
one encounters in the duality theory for schemes.
To motivate the construction we mention two additional properties:
\begin{enumerate}
\item $\varphi^!$ sends a dualizing complex for $R$ (if it exists)
to a dualizing complex for $A$, and
\item $\omega_{A/R}^\bullet = \varphi^!(R)$ is a kind of
relative dualizing complex: it lies in $D^b_{\textit{Coh}}(A)$ and restricts
to a dualizing complex on the fibres provided $R \to A$ is flat.
\end{enumerate}
These statements are Lemmas \href{https://stacks.math.columbia.edu/tag/0BZL}{lemma shriek dualizing algebraic (Tag 0BZL)} and
\href{https://stacks.math.columbia.edu/tag/0BZW}{lemma relative dualizing algebraic (Tag 0BZW)}.

\medskip\noindent
Let $\varphi : R \to A$ be a finite type homomorphism of Noetherian rings.
We will define a functor $\varphi^! : D(R) \to D(A)$ in the following way
\begin{enumerate}
\item If $\varphi : R \to A$ is surjective we set
$\varphi^!(K) = R\Hom(A, K)$. Here we use the functor
$R\Hom(A, -) : D(R) \to D(A)$ of
Section \href{https://stacks.math.columbia.edu/tag/0A6Z}{section trivial (Tag 0A6Z)}, and
\item in general we choose a surjection $\psi : P \to A$ with
$P = R[x_1, \ldots, x_n]$ and we set
$\varphi^!(K) = \psi^!(K \otimes_R^\mathbf{L} P)[n]$.
Here we use the functor
$- \otimes_R^\mathbf{L} P : D(R) \to D(P)$
of More on Algebra, Section \href{https://stacks.math.columbia.edu/tag/06Y5}{section derived base change (Tag 06Y5)}.
\end{enumerate}
Note the shift $[n]$ by the number of variables in the polynomial
ring. This construction is {\bf not} canonical and the functor
$\varphi^!$ will only be well defined up to a (nonunique) isomorphism of
functors\footnote{It is possible to make the construction canonical:
use $\Omega^n_{P/R}[n]$ instead of $P[n]$ in the
construction and use this in Lemma \href{https://stacks.math.columbia.edu/tag/0BZJ}{lemma well defined (Tag 0BZJ)}.
The material in this section becomes a lot more involved
if one wants to do this.}.



\begin{lemma}
\label{lemma-relative-dualizing-if-have-omega}
Let $\varphi : A \to B$ be a finite type homomorphism of Noetherian rings.
Let $\omega_A^\bullet$ be a dualizing complex for $A$. Set
$\omega_B^\bullet = \varphi^!(\omega_A^\bullet)$. Denote
$D_A(K) = R\Hom_A(K, \omega_A^\bullet)$ for $K \in D_{\textit{Coh}}(A)$
and
$D_B(L) = R\Hom_B(L, \omega_B^\bullet)$ for $L \in D_{\textit{Coh}}(B)$.
Then there is a functorial isomorphism
$$
\varphi^!(K) = D_B(D_A(K) \otimes_A^\mathbf{L} B)
$$
for $K \in D_{\textit{Coh}}(A)$.
\end{lemma}

\begin{proof}
Observe that $\omega_B^\bullet$ is a dualizing complex for $B$ by
Lemma \href{https://stacks.math.columbia.edu/tag/0BZL}{lemma shriek dualizing algebraic (Tag 0BZL)}.
Let $A \to B \to C$ be finite type homomorphisms of Noetherian rings.
If the lemma holds for $A \to B$ and $B \to C$, then the lemma holds for
$A \to C$. This follows from
Lemma \href{https://stacks.math.columbia.edu/tag/0BZT}{lemma composition shriek algebraic (Tag 0BZT)}
and the fact that $D_B \circ D_B \cong \text{id}$ by
Lemma \href{https://stacks.math.columbia.edu/tag/0A7C}{lemma dualizing (Tag 0A7C)}.
Thus it suffices to prove the lemma in case $A \to B$ is
a surjection and in the case where $B$ is a
polynomial ring over $A$.

\medskip\noindent
Assume $B = A[x_1, \ldots, x_n]$. Since $D_A \circ D_A \cong \text{id}$,
it suffices to prove
$D_B(K \otimes_A B) \cong D_A(K) \otimes_A B[n]$ for $K$
in $D_{\textit{Coh}}(A)$.
Choose a bounded complex $I^\bullet$ of injectives representing
$\omega_A^\bullet$. Choose a quasi-isomorphism
$I^\bullet \otimes_A B \to J^\bullet$ where $J^\bullet$
is a bounded complex of $B$-modules. Given a complex
$K^\bullet$ of $A$-modules, consider the obvious
map of complexes
$$
\Hom^\bullet(K^\bullet, I^\bullet) \otimes_A B[n]
\longrightarrow
\Hom^\bullet(K^\bullet \otimes_A B, J^\bullet[n])
$$
The left hand side represents $D_A(K) \otimes_A B[n]$ and the right hand
side represents $D_B(K \otimes_A B)$. Thus it suffices to prove this
map is a quasi-isomorphism if the cohomology modules
of $K^\bullet$ are finite $A$-modules. Observe that the
cohomology of the complex in degree $r$ (on either side)
only depends on finitely many of the $K^i$. Thus we may
replace $K^\bullet$ by a truncation, i.e., we may assume
$K^\bullet$ represents an object of $D^-_{\textit{Coh}}(A)$.
Then $K^\bullet$ is quasi-isomorphic to a bounded
above complex of finite free $A$-modules.
Therefore we may assume $K^\bullet$ is a bounded
above complex of finite free $A$-modules.
In this case it is easy to that the
displayed map is an isomorphism of complexes which finishes
the proof in this case.

\medskip\noindent
Assume that $A \to B$ is surjective. Denote $i_* : D(B) \to D(A)$
the restriction functor and recall that $\varphi^!(-) = R\Hom(A, -)$
is a right adjoint to $i_*$ (Lemma \href{https://stacks.math.columbia.edu/tag/0A70}{lemma right adjoint (Tag 0A70)}).
For $F \in D(B)$ we have
\begin{align*}
\Hom_B(F, D_B(D_A(K) \otimes_A^\mathbf{L} B))
& =
\Hom_B((D_A(K) \otimes_A^\mathbf{L} B) \otimes_B^\mathbf{L} F,
\omega_B^\bullet) \\
& =
\Hom_A(D_A(K) \otimes_A^\mathbf{L} i_*F, \omega_A^\bullet) \\
& =
\Hom_A(i_*F, D_A(D_A(K))) \\
& =
\Hom_A(i_*F, K) \\
& =
\Hom_B(F, \varphi^!(K))
\end{align*}
The first equality follows from More on Algebra, Lemma
\href{https://stacks.math.columbia.edu/tag/0A65}{lemma internal hom (Tag 0A65)} and the definition
of $D_B$. The second equality by the adjointness mentioned
above and the equality
$i_*((D_A(K) \otimes_A^\mathbf{L} B) \otimes_B^\mathbf{L} F) =
D_A(K) \otimes_A^\mathbf{L} i_*F$
(More on Algebra, Lemma \href{https://stacks.math.columbia.edu/tag/06Y6}{lemma derived base change (Tag 06Y6)}).
The third equality follows from More on Algebra, Lemma
\href{https://stacks.math.columbia.edu/tag/0A65}{lemma internal hom (Tag 0A65)}. The fourth because
$D_A \circ D_A = \text{id}$. The final equality by adjointness again.
Thus the result holds by the Yoneda lemma.
\end{proof}
\end{document}
